\documentclass[11pt,a4paper]{article}
\usepackage[T1]{fontenc}
\usepackage[utf8]{inputenc}
\usepackage{lmodern}
\usepackage{microtype}
\usepackage[a4paper,margin=27mm]{geometry}
\usepackage{amsmath,amssymb,amsthm,mathtools}
\usepackage{xcolor}
\usepackage{booktabs}
\usepackage{enumitem}
\usepackage{listings}
\usepackage[hidelinks]{hyperref}
\hypersetup{
  pdftitle={Exact Two-Dimensional SU(2) Yang-Mills in Lean},
  pdfauthor={Lluis Eriksson},
  pdfsubject={Weyl integration, heat-kernel convolution, Migdal invariance, and exact area law},
  pdfkeywords={Lean 4, Mathlib, SU(2), two-dimensional Yang-Mills, heat kernel, Migdal}
}

\definecolor{leanblue}{RGB}{36,83,125}
\definecolor{codegray}{RGB}{245,247,249}
\lstdefinelanguage{Lean}{
  morekeywords={theorem,lemma,def,noncomputable,section,namespace,open,by,fun,
    if,then,else,where,let,in,have,show,exact,simpa,rw,calc},
  sensitive=true,
  morecomment=[l]{--},
  morestring=[b]"
}
\lstset{
  language=Lean,
  basicstyle=\ttfamily\footnotesize,
  keywordstyle=\color{leanblue}\bfseries,
  commentstyle=\color{gray},
  backgroundcolor=\color{codegray},
  frame=single,
  rulecolor=\color{black!15},
  breaklines=true,
  columns=fullflexible,
  keepspaces=true,
  xleftmargin=0.5em,
  xrightmargin=0.5em
}

\newtheorem{theorem}{Theorem}[section]
\newtheorem{corollary}[theorem]{Corollary}
\newtheorem{proposition}[theorem]{Proposition}
\theoremstyle{remark}
\newtheorem{remark}[theorem]{Remark}

\title{\bfseries Exact Two-Dimensional $SU(2)$ Yang--Mills in Lean:\\
Weyl Integration, Heat-Kernel Convolution, Migdal Invariance,\\
and the Exact Simple-Loop Area Law}
\author{Lluis Eriksson}
\date{13 July 2026}

\begin{document}
\maketitle

\begin{abstract}
We give a kernel-checked Lean~4/Mathlib formalization of the exact
heat-kernel solution of two-dimensional $SU(2)$ Yang--Mills theory along the
chain from Haar measure to a concrete Migdal subdivision move.  The development
identifies the transported Haar probability on $SU(2)\simeq S^3$ with the
canonical spherical probability, proves the required all-order orbital
integration law, derives translated character convolution for every pair of
irreducible labels, passes from finite character sums to the infinite
heat-kernel semigroup by dominated convergence, and integrates an actual edge
shared by two faces.  We then define finite oriented disk cellulations
independently of any reduction tree, using vertices, paired half-edges, cyclic
face words, incidence, Euler characteristic, and positive face areas.  Every
connected dual graph admits a certified elimination schedule, every valid
schedule reduces to the heat kernel at total area, and any two choices give
the same amplitude.  For the first cyclic unreduced lattice, a three-face disk,
we additionally construct a triangular gauge-fixing equivalence that preserves
the original triple product Haar measure and prove equality with the reduced
amplitude.  We then close the global edge model: for every primal-connected,
facially well-formed finite disk cellulation, a rooted spanning tree selects
$V-1$ distinct physical edges and induces an explicit measurable equivalence
$SU(2)^E\simeq SU(2)^{V\setminus\{r\}}\times SU(2)^{E\setminus T}$ preserving
product Haar.  All facial holonomies transform simultaneously by basepoint
conjugation, the heat-kernel density loses every tree/gauge coordinate, and
the unreduced edge integral equals the gauge-fixed chord integral exactly.
The schedule is hidden from the public loop object,
which supplies a nontrivial exact-area-law package.  Combined with all-label
coefficient extraction, this yields the exact normalized simple-loop identity
\[
  \mathbb E[W_n(C)]
  =\exp\!\left[-\frac{n(n+2)}4\operatorname{Area}(C)\right].
\]
All displayed endpoints compile without \texttt{sorry}, \texttt{admit}, or
project-local axioms.  The mathematical identities are classical; the claimed
contribution is their concrete end-to-end machine-checked composition, not a
new analytic solution of two-dimensional Yang--Mills theory.
\end{abstract}

\section{Statement and claim boundary}

Let $dg$ be normalized Haar probability on $SU(2)$, let $\chi_n$ be the
character of the irreducible representation of dimension $d_n=n+1$, and put
\[
  c_n=\frac{n(n+2)}4,
  \qquad
  K_t(g)=\sum_{n\ge0}d_n e^{-tc_n}\chi_n(g),\qquad t>0.
\]
With convolution
\[
  (f*h)(g)=\int_{SU(2)}f(x)h(x^{-1}g)\,dx,
\]
the formalized chain is
\[
  \begin{aligned}
  \text{Haar on }S^3
  &\Longrightarrow \text{all-order orbital law}
  \Longrightarrow \chi_n*\chi_m\\
  &\Longrightarrow K_s*K_t=K_{s+t}
  \Longrightarrow \text{Migdal}
  \Longrightarrow \text{simple-loop coefficient}.
  \end{aligned}
\]

This result is distinct from the programme's earlier strong-coupling area-law
theorem~\cite{ErikssonStrongCoupling}.  The earlier theorem is a volume-uniform \emph{inequality} for finite
lattices and $SU(N_c)$ in an explicit convergence window, obtained from a
Koteck\'y--Preiss cluster expansion.  The present theorem is an \emph{exact
identity} in the two-dimensional $SU(2)$ heat-kernel model, obtained by Haar
integration and subdivision.  The former is cited context; it is not used to
discharge any exact two-dimensional statement here.

No claim is made here about four-dimensional Yang--Mills, a continuum
construction in four dimensions, or a new proof of the classical exact
solution.  A priority claim about formalization should remain conditional on a
systematic literature and repository search.

\section{From spherical Haar measure to the orbital law}

Write an element of $SU(2)$ in its standard form
\[
  g=\begin{pmatrix}a&b\\-\overline b&\overline a\end{pmatrix},
  \qquad |a|^2+|b|^2=1.
\]
Earlier modules construct a homeomorphism $SU(2)\simeq S^3$, realize left
multiplication as an ambient real-linear isometry, prove preservation of the
canonical \texttt{Measure.toSphere} measure, and use uniqueness of normalized
Haar measure to obtain literal equality of the two probability measures.

The new all-order step studies $u=|a|^2$.  In dimension four, the squared mass
in one complex rail is uniform on $[0,1]$.  Rather than assume this distribution,
the formal proof establishes its moment recurrence, proves
\[
  \int_{SU(2)}u^k\,dg=\frac1{k+1}\qquad(k\ge0),
\]
and then identifies measures on the compact interval using polynomial density.
The resulting literal endpoint is:

\begin{lstlisting}
theorem su2FirstRailMassMeasure_eq_uniform :
  su2FirstRailMassMeasure = su2UnitIntervalMeasure
\end{lstlisting}

For $x,y\in[-1,1]$, choose representatives with half-traces $x,y$ and average
their product over conjugation.  The preceding uniform law reduces the orbital
integral to a one-variable Chebyshev integral.  For the Chebyshev polynomial
$U_m$ of the second kind, the formalized formula is
\[
 \int_{SU(2)}
 U_m\!\left(\tfrac12\operatorname{tr}(kak^{-1}b)\right)dk
 =\frac{U_m(\tfrac12\operatorname{tr}a)
          U_m(\tfrac12\operatorname{tr}b)}{m+1}.
\]
The proof includes the endpoint cases $x=\pm1$; no division by a vanishing sine
is hidden in the statement.

\section{All-order translated character convolution}

Characters are implemented concretely as
\[
  \chi_n(g)=U_n\!\left(\tfrac12\operatorname{tr}g\right).
\]
Conjugacy invariance, left and right Haar invariance, Fubini, and the orbital
law give the translated orthogonality relation
\begin{equation}\label{eq:translated}
  \int_{SU(2)}\chi_n(x)\chi_m(x^{-1}g)\,dx
  =\delta_{nm}\frac{\chi_n(g)}{n+1}.
\end{equation}
This is stronger than orthogonality at the identity and is exactly the consumer
needed by the heat-kernel product.

\begin{theorem}[Machine-checked character convolution]
For every $n,m\in\mathbb N$ and $g\in SU(2)$,
\[
  (\chi_n*\chi_m)(g)=
  \begin{cases}
    \chi_n(g)/(n+1),&n=m,\\
    0,&n\ne m.
  \end{cases}
\]
\end{theorem}

The public Lean endpoints are listed below.  They are theorems about the actual
normalized Haar integral, not fields projected from an abstract character
package.
\begin{lstlisting}
integral_su2CharacterReal_mul_translate
su2CharacterChebyshev_convolution
\end{lstlisting}

\section{Infinite heat-kernel semigroup}

For a cutoff $N$, expand both finite kernels, interchange the finite sums and
the Haar integral, and apply \eqref{eq:translated}.  The Kronecker selector
leaves a single sum and the Casimir weights multiply:
\[
  e^{-s c_n}e^{-t c_n}=e^{-(s+t)c_n}.
\]
This proves the finite convolution identity.  The infinite theorem then uses
the already established uniform Casimir majorant for positive times.  Partial
kernels converge pointwise and uniformly; the product is bounded by the
product of two summable majorants, so dominated convergence passes the Haar
integral to the limit.

\begin{theorem}[Concrete heat semigroup]
For $s,t>0$ and every $g\in SU(2)$,
\[
  \int_{SU(2)}K_s(x)K_t(x^{-1}g)\,dx=K_{s+t}(g).
\]
Equivalently, as functions, $K_s*K_t=K_{s+t}$.
\end{theorem}

The corresponding Lean declarations, including the function-level consumer,
are:
\begin{lstlisting}
su2HeatKernelPartial_convolution
su2HeatKernel_convolution
su2HeatKernel_convolution_consumer
\end{lstlisting}

\section{A concrete Migdal move}

Consider two adjacent faces of areas $s$ and $t$.  Let $a,b$ encode the
external boundary holonomies and integrate their shared oriented edge $x$.
The actual two-face density in the development is
\[
  D_{s,t}^{a,b}(x)=K_s(ax)K_t(x^{-1}b).
\]
Left invariance changes variables $y=ax$ and the heat semigroup gives
\begin{equation}\label{eq:migdal}
  \int_{SU(2)}D_{s,t}^{a,b}(x)\,dx
  =K_{s+t}(ab).
\end{equation}

\begin{theorem}[Two-face Migdal subdivision]
Equation \eqref{eq:migdal} holds for all $a,b\in SU(2)$ and all $s,t>0$.
It is proved by a genuine Haar integration over the shared edge and invokes
the concrete infinite heat-kernel semigroup.
\end{theorem}

This closes the first nontrivial subdivision move: two holonomy variables, two
faces, and one integrated shared edge.  The endpoints below contain no package
hypothesis.
\begin{lstlisting}
su2Migdal_twoFace_merge
su2Migdal_subdivision_invariant
\end{lstlisting}

\section{A cyclic product-Haar gauge fixing}

The two-face move alone does not identify a reduced evaluator with the original
lattice integral when the face-dual graph has cycles.  We therefore formalize
the first cyclic test explicitly.  Consider a disk split into three bounded
faces by three spokes meeting at one interior vertex.  Its face dual is
$K_3$, equivalently the three-cycle.  Write $x,y,z\in SU(2)$ for the original
spoke-edge variables and $a,b,c\in SU(2)$ for the fixed boundary arcs.  The
three face holonomies are
\[
  xay^{-1},\qquad ybz^{-1},\qquad zcx^{-1}.
\]
Thus the unreduced density and measure are
\begin{equation}\label{eq:unreduced-three}
  K_{t_1}(xay^{-1})K_{t_2}(ybz^{-1})K_{t_3}(zcx^{-1})
  \,dx\,dy\,dz.
\end{equation}

Gauge transformations at the interior vertex act by simultaneous left
multiplication.  Lean constructs the measurable equivalence
\[
  \Phi(x,y,z)=(x,x^{-1}y,x^{-1}z),
  \qquad
  \Phi^{-1}(u,v,w)=(u,uv,uw).
\]
The skew-product theorem and left invariance of normalized Haar prove the
literal measure identity
\[
  \Phi_*\bigl(dx\,dy\,dz\bigr)=dx\,dy\,dz.
\]
After substitution, conjugation invariance of the heat kernel removes the
first coordinate and gives
\[
  K_{t_1}(av^{-1})K_{t_2}(vbw^{-1})K_{t_3}(wc).
\]
The pure-gauge integral has mass one.  Applying the reversed two-face Migdal
identity first to $w$ and then to $v$ yields the fully unreduced result
\begin{equation}\label{eq:three-spoke-result}
 \int_{SU(2)^3}
  K_{t_1}(xay^{-1})K_{t_2}(ybz^{-1})K_{t_3}(zcx^{-1})
  \,dx\,dy\,dz
  =K_{t_1+t_2+t_3}(abc).
\end{equation}

The proof does not store a Jacobian certificate or assume an abstract gauge
package.  The map, its inverse, measurability, product-Haar preservation,
density factorization, Fubini integrability, and both edge integrations are
all kernel checked.  The principal endpoints are:
\begin{lstlisting}
su2ThreeSpokeFaceDual_connected
su2ThreeSpokeGaugeFix_measurePreserving
su2ThreeSpoke_density_gaugeFactorization
su2ThreeSpoke_unreducedIntegral_eq_gaugeFixedIntegral
su2ThreeSpoke_unreducedIntegral_eq_heatKernel
\end{lstlisting}

Equation~\eqref{eq:three-spoke-result} is the first concrete product-Haar
instance.  It is now subsumed by the uniform original-edge theorem below,
whose spanning-tree gauge fixing is constructed directly from the incidence
data of the independent disk object.

\subsection{Uniform finite-valence theorem}

The preceding three-spoke map is not intrinsically three-dimensional.  Let
$I$ be any finite type and let
\[
  d\mu_I=\bigotimes_{i\in I}dg_i
\]
be the normalized product Haar probability on $SU(2)^I$.  The formalization
constructs, without choosing an enumeration of $I$, the measurable equivalence
\[
 \begin{aligned}
 \Phi_I&\colon SU(2)\times SU(2)^I\to SU(2)\times SU(2)^I,\\
 \Phi_I(x,y)&=\bigl(x,(x^{-1}y_i)_{i\in I}\bigr).
 \end{aligned}
\]
with inverse $(u,v)\mapsto(u,(uv_i)_i)$.  The product measure is Mathlib's
\texttt{Measure.pi}.  Using finite-product Haar invariance and the skew-product
theorem, Lean proves
\begin{equation}\label{eq:finite-valence-mp}
  (\Phi_I)_*(dg\otimes d\mu_I)=dg\otimes d\mu_I
\end{equation}
for every finite $I$.  It also proves that the diagonal action
$(x,y_i)\mapsto(ux,uy_i)$ preserves the same measure.

The result is consumed at the integral level.  If
$F:SU(2)\times SU(2)^I\to\mathbb C$ is invariant under that diagonal action,
then
\begin{equation}\label{eq:identity-slice}
 \int F(x,y)\,dg(x)\,d\mu_I(y)
 =\int F(1,y)\,d\mu_I(y).
\end{equation}
More generally, the source contains an explicit factorization theorem for any
density whose pullback through $\Phi_I^{-1}$ loses its first coordinate.  The
principal uniform endpoints are:
\begin{lstlisting}
su2FiniteStarDiagonalLeft_measurePreserving
su2FiniteStarGaugeFix_measurePreserving
su2FiniteStar_integral_eq_gaugeFixed
su2FiniteStar_integral_eq_identitySlice
\end{lstlisting}

Equations~\eqref{eq:finite-valence-mp}--\eqref{eq:identity-slice} remove the
fixed-valence limitation of the local gauge step.  The next two subsections
compose them with an ordered spanning tree, first for vertex coordinates and
then for all original edge variables and all facial holonomies simultaneously.

\subsection{Global rooted-tree vertex coordinates}

The development now performs the global composition at the level of vertex
coordinates.  A value of \texttt{SU2RootedTreeOrder n} records a rooted tree on
$n+1$ vertices by adjoining one leaf at a time, together with the already
present parent of each new leaf.  Its parent--child list is embedded in the
final index type.  For every finite connected simple graph, Lean constructs an
equivalence
\[
  e:\operatorname{Fin}(n+1)\simeq V
\]
and such an order for which every recorded parent--child pair is an actual
graph edge.  The proof removes a vertex whose complement remains connected,
recurses on the induced graph, and reattaches the vertex across a witnessed
edge.

For vertex variables $x_0,\ldots,x_n\in SU(2)$, the associated single global
map retains the root and replaces each later variable by its tree increment:
\begin{equation}\label{eq:global-tree-coordinates}
  \Psi_T(x)_0=x_0,
  \qquad
  \Psi_T(x)_k=x_{p(k)}^{-1}x_k\quad(k>0).
\end{equation}
It is implemented recursively by splitting the last coordinate, applying the
triangular Haar shear, transforming the older coordinates, and reassembling
the finite product.  The inverse is part of the measurable equivalence, not an
existence assertion.  Induction and finite-product Haar invariance prove
\begin{equation}\label{eq:global-tree-haar}
  (\Psi_T)_*\!\left(\bigotimes_{v\in V}dg_v\right)
  =\bigotimes_{v\in V}dg_v.
\end{equation}

For a \texttt{SU2FiniteDiskCellulation}, primal adjacency is derived from an
actual oriented half-edge and made into a simple graph.  Under the explicit
hypothesis that this primal graph is connected, Lean chooses the certified
tree internally, transports \eqref{eq:global-tree-coordinates} through the
vertex enumeration, and returns a global measure-preserving equivalence on
\texttt{C.Vertex -> SU2}.  The principal endpoints are:
\begin{lstlisting}
SU2RootedTreeOrder.exists_validForGraph_of_connected
SU2RootedTreeOrder.coordinateEquiv_measurePreserving
SU2FiniteDiskCellulation.exists_rootedSpanningTree_of_primal_connected
SU2FiniteDiskCellulation.RootedSpanningTree.vertexCoordinateEquiv_measurePreserving
SU2FiniteDiskCellulation.exists_rootedTreeVertexCoordinateEquiv_measurePreserving
\end{lstlisting}

The edge module consumes this tree and proves the stronger decomposition
\[
 SU(2)^E\simeq SU(2)^{V\setminus\{r\}}\times SU(2)^{E\setminus T},
\]
as a literal measurable equivalence.  Each construction step chooses an
oriented parent--child half-edge; strict ordering of parents before children
proves that the resulting physical tree edges are injective.  The complement
$E\setminus T$ is therefore an actual finite subtype, not a cardinality-only
placeholder.  After correcting the stored orientation of each tree edge, the
map reconstructs potentials $g_r=1$ and
$g_k=g_{p(k)}a_k$, and sends each chord to
\[
  X_e=g_{s(e)}U_e g_{t(e)}^{-1}.
\]
Finite-product relabelling, coordinatewise inversion and a Haar-preserving
skew product prove preservation of normalized product Haar.

The physical edge record \texttt{SU2EdgeConnectedDiskCellulation} additionally
requires an explicit once-around, duplicate-free boundary enumeration for
each bounded face.  This condition is stated because the older combinatorial
record did not exclude a formally declared face with no incident half-edge.
For every well-formed face $f$, Lean proves simultaneously
\[
  H_f(g\mathbin{\cdot}U)
  =g_{v_f}H_f(U)g_{v_f}^{-1}.
\]
Consequently the complete heat-kernel density is gauge invariant.  Under the
inverse edge equivalence it factors only through the chord block, and the
normalized gauge coordinates integrate to one:
\begin{equation}\label{eq:unreduced-global-edge}
 \int_{SU(2)^E}\prod_f K_{t_f}(H_f(U))\,dU
 =\int_{SU(2)^{E\setminus T}}
   \prod_f K_{t_f}(H_f(U_T=1,X))\,dX.
\end{equation}
Both sides of \eqref{eq:unreduced-global-edge} integrate every physical edge,
including the exterior boundary edges, and are therefore scalars.  They must
not be identified with the two-boundary amplitude $A_C(a,b)$ below, which
retains an exterior holonomy and equals $K_{\sum_f t_f}(ab)$.  A pointwise
edge-model derivation of that amplitude requires a separately defined
boundary-conditioned edge integral and remains a distinct bridge theorem.
The principal new endpoints are:
\begin{lstlisting}
SU2FiniteDiskCellulation.RootedSpanningTree.treeEdge_injective
SU2FiniteDiskCellulation.RootedSpanningTree.globalEdgeGaugeEquiv_measurePreserving
SU2FiniteDiskCellulation.RootedSpanningTree.gaugeTransform_eq_gaugeFixedEdgeConfiguration
SU2EdgeConnectedDiskCellulation.allFaceHolonomies_gaugeTransform
SU2EdgeConnectedDiskCellulation.edgeHeatKernelDensity_globalEdgeFactorization
SU2EdgeConnectedDiskCellulation.unreducedEdgeIntegral_eq_chordGaugeFixedIntegral
\end{lstlisting}

\section{Independent finite cellulations and elimination}

The analytic evaluator used by the development is a labelled binary schedule.
If a schedule node has left and right subproblems $L,R$, its amplitude is
defined recursively by
\[
  A_C(a,b)=\int_{SU(2)} A_L(a,x)A_R(x^{-1},b)\,dx.
\]
The earlier tree module proves by structural induction that
\[
  A_C(a,b)=K_{T(C)}(ab).
\]
At each node the induction hypotheses expose two heat kernels and the proof
invokes the concrete two-face theorem.

The new point is that the geometric input is no longer this evaluator syntax.
The new independent cellulation structure stores finite types of
vertices, edges, half-edges, and bounded faces; a fixed-point-free involution
pairs the two half-edges of every edge; source and target incidence is linked
to a cyclic face-successor permutation; \texttt{none} denotes the exterior
face; and the disk Euler relation $V-E+F=1$ is part of the certificate.  The
dual face graph is derived from the two sides of internal edges.

A labelled elimination schedule is valid when its leaves list every bounded
face exactly once and each binary merge is witnessed by a genuine adjacency
in the derived dual graph.  Lean proves the following existence theorem for an
arbitrary finite connected dual graph.  If the graph has one vertex, the
schedule is a leaf.  Otherwise, a vertex whose deletion leaves a connected
induced subgraph is removed, the proof recurses on the smaller finite type,
and the vertex is reattached across an actual adjacent edge.  Thus no
acyclicity assumption is used.

For any valid schedule $S$, permutation of the face list gives
$T(S)=\sum_f t_f$.  The tree reduction then proves
\begin{equation}\label{eq:global-reduction}
  A_{C,S}(a,b)=K_{\sum_f t_f}(ab).
\end{equation}
Consequently, for all valid $S_1,S_2$,
\[
  A_{C,S_1}(a,b)=A_{C,S_2}(a,b).
\]
The public structure \texttt{SU2ConnectedDiskCellulation} stores only the
cellulation and connectedness of its dual.  A schedule is chosen internally
from the existence theorem; the preceding equality proves that every other
valid choice is observationally identical.  Its principal declarations are:
\begin{lstlisting}
SU2FaceEliminationSchedule.exists_schedule_of_connected_graph
SU2FaceEliminationSchedule.exists_valid_schedule_of_dual_connected
SU2FaceEliminationSchedule.amplitude_eq_of_valid_schedules
SU2ConnectedDiskCellulation.amplitude_eq_heatKernel
SU2ConnectedDiskCellulation.amplitude_eq_any_valid_schedule
\end{lstlisting}

For each representation label $n$, the formalization defines a plane-loop
theory whose loops are connected independently specified cellulations, whose
area is $\sum_f t_f$, whose string tension is $n(n+2)/4$, and whose Wilson
functional integrates the schedule-independent amplitude.  Equation
\eqref{eq:global-reduction} followed by
coefficient extraction gives a nontrivial concrete instance
\texttt{su2ConnectedDiskExactAreaLawPackage}.  The generic public interface is
then consumed by
\texttt{su2ConnectedDisk\_simpleLoop\_areaLaw\_exact}; no field of that
package remains assumed.  The trivial-label case also proves the normalization
\[
  \int_{SU(2)} A_C(1,g)\,dg=1.
\]

For a general dual graph with cycles, the schedule theorem concerns the
explicitly defined post-gauge-fixed amplitude.  Equation
\eqref{eq:unreduced-global-edge} now supplies the uniform original-edge gauge
reduction for every facially well-formed, primal- and dual-connected disk
cellulation; the earlier three-cycle computation remains a concrete diagnostic
instance of the general construction.  It does not by itself identify the
fully boundary-integrated chord scalar with the boundary-dependent schedule
amplitude.

\section{Exact simple-loop coefficient}

The same repository proves, for every label $n$ and positive area $t$, that
the normalized Wilson character $W_n=\chi_n/(n+1)$ extracts the corresponding
heat coefficient:
\[
  \int_{SU(2)}W_n(g)K_t(g)\,dg
  =e^{-t c_n}
  =\exp\!\left[-t\frac{n(n+2)}4\right].
\]
The proof first establishes the identity for partial character sums and then
passes to the full kernel by dominated convergence.  Its final consumer is:

\begin{samepage}
\begin{lstlisting}
su2_exact_simpleLoop_areaLaw
\end{lstlisting}
\end{samepage}

Together with the iterated Migdal reduction, this is the exact
two-dimensional route for connected finite combinatorial disks after the
certified elimination step: the choice of schedule does not change the
density, and the effective face has the Casimir coefficient prescribed by its
total area.

\section{Audit and reproducibility}

The checked source is built with Lean~4 and Mathlib.  The public root module
imports the concrete orbit, character-convolution, heat-semigroup, and exact
area-law modules.  A repository-wide placeholder audit finds no occurrence of
\texttt{sorry}, \texttt{admit}, or a project-local \texttt{axiom} declaration
in the Lean development.  Printing axioms for the fifty-two audited public endpoints
returns only the standard kernel dependencies
\[
  \texttt{propext},\qquad \texttt{Classical.choice},\qquad
  \texttt{Quot.sound}.
\]
The toolchain and Mathlib revision are pinned in the repository.  GitHub
Actions builds the public root target from a clean checkout and separately
rejects placeholders.

\paragraph{Authoritative proof artifact.}
The complete machine-checked proofs are contained in the public repository
\texttt{lluiseriksson/lean-2d-yang-mills}~\cite{FormalArtifact}.  The immutable
audited snapshot for this manuscript is commit
\texttt{52e0e6f54e40a891c9dc092f730daba26f100bcd}.  In particular,
\texttt{SU2RootedTreeGaugeFixing.lean} contains the connected-graph spanning-tree
construction and global product-Haar theorem; the preceding modules contain
the Haar--sphere identification, orbital integration, character convolution,
heat semigroup, Migdal move, finite-cellulation reduction, and exact coefficient
extraction.  \texttt{AuditMigdalClosure.lean} prints the kernel dependencies of
all fifty-two public endpoints.  Thus the repository is not supplementary
numerics or pseudocode: it is the compilable proof object to which the theorem
claims refer.

\medskip\noindent\textbf{Principal new declarations.}\par\smallskip
\begin{lstlisting}
su2FirstRailMoment_recurrence
su2FirstRailMoment_eq
su2FirstRailMassMeasure_eq_uniform
intervalIntegral_chebyshevU_product_formula
integral_su2FirstEntryMass_eq_unitInterval
integral_su2Haar_orbit_character_general
integral_su2CharacterReal_mul_translate
su2CharacterChebyshev_convolution
su2HeatKernelPartial_convolution
su2HeatKernel_convolution
su2Migdal_twoFace_merge
su2Migdal_subdivision_invariant
su2_exact_simpleLoop_areaLaw
SU2PlanarFaceTree.internalEdgeCount_eq_faceCount_sub_one
su2PlanarCellulationAmplitude_eq_heatKernel
su2PlanarCellulationAmplitude_eq_effectiveFace
su2TreePlanar_wilsonExpectation_eq_casimir
su2TreePlanarExactAreaLawPackage
su2TreePlanar_simpleLoop_areaLaw_exact
integral_su2PlanarCellulationAmplitude_eq_one
SU2FaceEliminationSchedule.exists_schedule_of_connected_graph
SU2FaceEliminationSchedule.exists_valid_schedule_of_dual_connected
SU2FaceEliminationSchedule.amplitude_eq_of_valid_schedules
SU2ConnectedDiskCellulation.amplitude_eq_heatKernel
SU2ConnectedDiskCellulation.amplitude_eq_any_valid_schedule
su2ConnectedDisk_wilsonExpectation_eq_casimir
su2ConnectedDiskExactAreaLawPackage
su2ConnectedDisk_simpleLoop_areaLaw_exact
integral_su2ConnectedDiskAmplitude_eq_one
su2ThreeSpokeGaugeFix_measurePreserving
su2ThreeSpoke_density_gaugeFactorization
su2ThreeSpoke_gaugeFixedIntegral_eq_heatKernel
su2ThreeSpoke_unreducedIntegral_eq_gaugeFixedIntegral
su2ThreeSpoke_unreducedIntegral_eq_heatKernel
su2FiniteStarDiagonalLeft_measurePreserving
su2FiniteStarGaugeFix_measurePreserving
su2FiniteStar_integral_eq_gaugeFixed
su2FiniteStar_integral_eq_identitySlice
SU2RootedTreeOrder.exists_validForGraph_of_connected
SU2RootedTreeOrder.coordinateEquiv_measurePreserving
SU2FiniteDiskCellulation.exists_rootedSpanningTree_of_primal_connected
SU2FiniteDiskCellulation.RootedSpanningTree.vertexCoordinateEquiv_measurePreserving
SU2FiniteDiskCellulation.exists_rootedTreeVertexCoordinateEquiv_measurePreserving
SU2FiniteDiskCellulation.RootedSpanningTree.globalEdgeGaugeEquiv_measurePreserving
SU2FiniteDiskCellulation.RootedSpanningTree.gaugeTransform_eq_gaugeFixedEdgeConfiguration
SU2EdgeConnectedDiskCellulation.allFaceHolonomies_gaugeTransform
SU2EdgeConnectedDiskCellulation.unreducedEdgeIntegral_eq_chordGaugeFixedIntegral
\end{lstlisting}

\section{Scope and next frontier}

The present artifact closes the compact-group analytic chain, introduces an
independent finite oriented disk-cellulation object, constructs elimination
schedules for every connected dual graph, and proves independence of the
schedule for the reduced amplitude.  Dual cycles are allowed at this level,
and the full original-edge model is linked to its gauge-fixed chord integral by
a global lattice gauge-fixing equivalence.  The construction is uniform in the
finite cellulation, proves simultaneous facial covariance and removes all
$V-1$ gauge coordinates exactly.

The artifact does not identify the abstract combinatorial maps with planar
isotopy classes, treat nonsimple or intersecting loops, construct a continuum
measure, or evaluate positive-area higher-genus partition functions.  It also
does not yet identify a boundary-conditioned original-edge integral with the
schedule amplitude for every cellulation.  That pointwise boundary bridge,
followed by the geometric and continuum questions, forms the remaining
frontier.

\begin{thebibliography}{9}
\bibitem{Migdal}
A.~A. Migdal, \emph{Recursion equations in gauge field theories},
Sov. Phys. JETP \textbf{42} (1975), 413--418.

\bibitem{Witten}
E.~Witten, \emph{On quantum gauge theories in two dimensions},
Commun. Math. Phys. \textbf{141} (1991), 153--209.

\bibitem{Driver}
B.~K. Driver, \emph{Yang--Mills theory and the heat equation on compact Lie
groups}, J. Funct. Anal. \textbf{110} (1992), 1--34.

\bibitem{Lean}
L.~de Moura and S.~Ullrich, \emph{The Lean 4 theorem prover and programming
language}, CADE-28, Lecture Notes in Computer Science 12699 (2021), 625--635.

\bibitem{Mathlib}
The Mathlib Community, \emph{The Lean mathematical library},
CPP 2020, 367--381.

\bibitem{ErikssonStrongCoupling}
L.~Eriksson, \emph{A machine-checked volume-uniform Wilson-loop area law via a
formalized cluster expansion},
\href{https://ai.vixra.org/pdf/2607.0005v1.pdf}{ai.viXra:2607.0005} (2026).

\bibitem{FormalArtifact}
L.~Eriksson, \emph{lean-2d-yang-mills: machine-checked exact two-dimensional
$SU(2)$ Yang--Mills}, GitHub repository,
\href{https://github.com/lluiseriksson/lean-2d-yang-mills/tree/52e0e6f54e40a891c9dc092f730daba26f100bcd}{audited Lean snapshot 52e0e6f}
(2026).
\end{thebibliography}

\end{document}
